\clearpage
\SourceLeaf{106}\pdfbookmark[1]{Sectio tertia.}{sectio-tertia}
\begin{center}
{\large SECTIO TERTIA.}\par\medskip
\emph{Relationes inter locos plures in orbita.}\par\bigskip
\hypertarget{native-106-article078}{}\pdfbookmark[2]{78.}{article-078}78.
\end{center}
Comparatio duorum pluriumve locorum corporis coelestis tum in orbita tum in spatio tantam propositionum elegantium copiam subministrat, ut volumen integrum facile complerent. Nostrum vero propositum non eo tendit, ut hoc argumentum fertile exhauriamus, sed eo potissimum, ut amplum apparatum subsidiorum ad solutionem problematis magni de determinatione orbitarum incognitarum ex observationibus inde adstruamus: quamobrem neglectis quae ab instituto nostro nimis aliena essent, eo diligentius omnia quae ullo modo illuc conducere possunt evolvemus. Disquisitionibus ipsis quasdam propositiones trigonometricas praemittimus, ad quas, quum frequentioris usus sint, saepius recurrere oportet.
\par
I. Denotantibus $A$, $B$, $C$ angulos quoscunque, habetur
\[
\begin{aligned}
\sin A\sin(C-B)+\sin B\sin(A-C)+\sin C\sin(B-A)&=0\\
\cos A\sin(C-B)+\cos B\sin(A-C)+\cos C\sin(B-A)&=0.
\end{aligned}
\]
\par
II. Si duae quantitates $p$, $P$ ex aequationibus talibus
\[
\begin{aligned}
p\sin(A-P)&=a\\
p\sin(B-P)&=b
\end{aligned}
\]
determinandae sunt, hoc fiet generaliter adiumento formularum
\SourceLeaf{107}
\[
\begin{aligned}
p\sin(B-A)\sin(H-P)&=b\sin(H-A)-a\sin(H-B)\\
p\sin(B-A)\cos(H-P)&=b\cos(H-A)-a\cos(H-B),
\end{aligned}
\]
in quibus $H$ est angulus arbitrarius. Hinc deducuntur (art. 14, II) angulus $H-P$ atque $p\sin(B-A)$; et hinc $P$ et $p$. Plerumque conditio adiecta esse solet, ut $p$ esse debeat quantitas positiva, unde ambiguitas in determinatione anguli $H-P$ per tangentem suam deciditur; deficiente autem illa conditione, ambiguitatem ad lubitum decidere licebit. Ut calculus commodissimus sit, angulum arbitrarium $H$ vel $=A$ vel $=B$ vel $=\tfrac{1}{2}(A+B)$ statuere conveniet. In casu priori aequationes ad determinandum $P$ et $p$ erunt
\[
\begin{aligned}
p\sin(A-P)&=a\\
p\cos(A-P)&=\frac{b-a\cos(B-A)}{\sin(B-A)}\mathpunct{\cdot}
\end{aligned}
\]
In casu secundo aequationes prorsus analogae erunt; in casu tertio autem
\[
\begin{aligned}
p\sin(\tfrac{1}{2}A+\tfrac{1}{2}B-P)&=\frac{b+a}{2\cos\tfrac{1}{2}(B-A)}\\
p\cos(\tfrac{1}{2}A+\tfrac{1}{2}B-P)&=\frac{b-a}{2\sin\tfrac{1}{2}(B-A)}\mathpunct{\cdot}
\end{aligned}
\]
\par
Quodsi itaque angulus auxiliaris $\zeta$ introducitur, cuius tangens $=\frac{a}{b}$, invenietur $P$ per formulam
\[
\operatorname{tang}(\tfrac{1}{2}A+\tfrac{1}{2}B-P)=\operatorname{tang}(45^0+\zeta)\operatorname{tang}\tfrac{1}{2}(B-A)
\]
ac dein $p$ per aliquam formularum praecedentium, ubi
\[
\begin{aligned}
\tfrac{1}{2}(b+a)&=\sin(45^0+\zeta)\BareRoot{\frac{ab}{\sin2\zeta}}\\
&=\frac{a\sin(45^0+\zeta)}{\sin\zeta\BareRoot{2}}=\frac{b\sin(45^0+\zeta)}{\cos\zeta\BareRoot{2}}\\[.4em]
\tfrac{1}{2}(b-a)&=\cos(45^0+\zeta)\BareRoot{\frac{ab}{\sin2\zeta}}\\
&=\frac{a\cos(45^0+\zeta)}{\sin\zeta\BareRoot{2}}=\frac{b\cos(45^0+\zeta)}{\cos\zeta\BareRoot{2}}\mathpunct{\cdot}
\end{aligned}
\]
\par
III. Si $p$ et $P$ determinandae sunt ex aequationibus
\[
\begin{aligned}
p\cos(A-P)&=a\\
p\cos(B-P)&=b,
\end{aligned}
\]
omnia in II exposita statim applicari possent, si modo illic pro $A$ et $B$ ubique scriberetur $90^0+A$, $90^0+B$: sed ut usus eo commodior sit, formulas evolutas apponere non piget. Formulae generales erunt
\SourceLeaf{108}
\[
\begin{aligned}
p\sin(B-A)\sin(H-P)&=-b\cos(H-A)+a\cos(H-B)\\
p\sin(B-A)\cos(H-P)&=b\sin(H-A)-a\sin(H-B).
\end{aligned}
\]
Transeunt itaque, pro $H=A$, in
\[
\begin{aligned}
p\sin(A-P)&=\frac{a\cos(B-A)-b}{\sin(B-A)}\\
p\cos(A-P)&=a.
\end{aligned}
\]
Pro $H=B$, formam similem obtinent; pro $H=\tfrac{1}{2}(A+B)$ autem fiunt
\[
\begin{aligned}
p\sin(\tfrac{1}{2}A+\tfrac{1}{2}B-P)&=\frac{a-b}{2\sin\tfrac{1}{2}(B-A)}\\
p\cos(\tfrac{1}{2}A+\tfrac{1}{2}B-P)&=\frac{a+b}{2\cos\tfrac{1}{2}(B-A)}\mathpunct{\raisebox{.45ex}{,}}
\end{aligned}
\]
ita ut introducto angulo auxiliari $\zeta$, cuius tangens $=\frac{a}{b}$, fiat
\[
\operatorname{cotang}(\tfrac{1}{2}A+\tfrac{1}{2}B-P)=\operatorname{cotang}(\zeta-45^0)\operatorname{tang}\tfrac{1}{2}(B-A).
\]
\par
Ceterum si $p$ immediate ex $a$ et $b$ sine praevio computo anguli $P$ determinare cupimus, habemus formulam
\[
p\sin(B-A)=\BareRoot{(aa+bb-2ab\cos(B-A))}
\]
tum in problemate praesente tum in II.
